% SPDX-License-Identifier: MIT
% Copyright (c) 2026 Evgeny Baulin
% Optimization Methods, Seminar 06: theory handout (English).
% Generated by checks/06. Second-Order Methods Part Two.ipynb from web/06/js/content-*.js and content/students.json. Do not edit by hand.
% Compile with: latexmk -xelatex -interaction=nonstopmode -halt-on-error Theory_en.tex
\documentclass[11pt,a4paper]{article}

\usepackage{fontspec}
\setmainfont{cmunrm.otf}[BoldFont=cmunbx.otf, ItalicFont=cmunti.otf, BoldItalicFont=cmunbi.otf]
\setsansfont{cmunss.otf}[BoldFont=cmunsx.otf, ItalicFont=cmunsi.otf, BoldItalicFont=cmunso.otf]
\setmonofont{cmuntt.otf}[BoldFont=cmuntb.otf, ItalicFont=cmunit.otf, BoldItalicFont=cmuntx.otf]
\usepackage{polyglossia}
\setdefaultlanguage{english}

\usepackage[a4paper,left=2.3cm,right=2.3cm,top=2.4cm,bottom=2.4cm,headheight=14pt]{geometry}
\usepackage{amsmath,amssymb}
\usepackage{graphicx}
\usepackage{array,booktabs,longtable}
\usepackage{enumitem}
\usepackage{xcolor}
\usepackage[most]{tcolorbox}
\usepackage{titlesec}
\usepackage{fancyhdr}
\usepackage{needspace}
\usepackage[expansion=false]{microtype}
\usepackage[colorlinks=true,linkcolor=hseblue,urlcolor=hseblue,citecolor=hseblue]{hyperref}

\definecolor{hseblue}{RGB}{16,45,105}
\definecolor{hseaccent}{RGB}{198,40,40}
\definecolor{hsegray}{RGB}{108,114,126}
\definecolor{hselight}{RGB}{240,243,248}
\definecolor{hsewarm}{RGB}{253,244,232}

\hypersetup{pdftitle={Optimization Methods. Seminar 06. Theory handout}, pdfauthor={Evgeny Baulin}}

\titleformat{\section}{\color{hseblue}\Large\bfseries}{\thesection}{0.6em}{}
\titleformat{\subsection}{\color{hseblue}\large\bfseries}{\thesubsection}{0.6em}{}
\titlespacing*{\section}{0pt}{16pt plus 4pt minus 2pt}{8pt plus 2pt}
\titlespacing*{\subsection}{0pt}{12pt plus 3pt minus 2pt}{6pt plus 2pt}
\pretocmd{\section}{\Needspace{7\baselineskip}}{}{\errmessage{cannot patch section}}
\pretocmd{\subsection}{\Needspace{5\baselineskip}}{}{\errmessage{cannot patch subsection}}

\tcbset{handoutbox/.style={enhanced, breakable, sharp corners, boxrule=0pt, left=7pt, right=7pt, top=5pt, bottom=5pt, before skip=8pt plus 2pt, after skip=8pt plus 2pt}}
\newtcolorbox{examplebox}{handoutbox, colback=hselight, colframe=hselight, before upper={\textcolor{hseblue}{\bfseries Example.}\ }}
\newtcolorbox{whybox}{handoutbox, colback=white, colframe=hseblue, leftrule=2.5pt, top=3pt, bottom=3pt, before upper={\textcolor{hseaccent}{\bfseries Why.}\ }}
\newtcolorbox{answerbox}{handoutbox, colback=hsewarm, colframe=hsewarm, before upper={\textcolor{hseaccent}{\bfseries Answer.}\ }}
\newtcolorbox{conventionbox}{handoutbox, colback=hselight, colframe=hseblue, leftrule=2.5pt, before upper={\textcolor{hseblue}{\bfseries Notation.}\par\smallskip}}

\newcommand{\handoutfigure}[3]{\par\medskip\noindent\begin{minipage}{\linewidth}\centering
  \includegraphics[width=#1\linewidth]{figures/#2.pdf}\par\smallskip
  \begin{minipage}{0.94\linewidth}\small\textcolor{hseblue}{\bfseries Figure.}\ #3\end{minipage}\end{minipage}\par\medskip}
\newcommand{\problemhead}[3]{\par\bigskip\noindent{\color{hseblue}\large\bfseries #1.\ #2}\par\nopagebreak
  \noindent{\small\itshape\color{hsegray}#3}\par\nopagebreak\smallskip}
\newtcolorbox{definitionbox}[1]{handoutbox, colback=white, colframe=hseblue, leftrule=2.5pt, before upper={\textcolor{hseblue}{\bfseries #1}\ }}
\newtcolorbox{theorembox}[1]{handoutbox, colback=hselight, colframe=hseblue, leftrule=2.5pt, before upper={\textcolor{hseblue}{\bfseries #1}\ }}
\newtcolorbox{proofbox}{handoutbox, colback=white, colframe=hsegray, leftrule=1.2pt, top=3pt, bottom=3pt, before upper={\textcolor{hsegray}{\bfseries\itshape Proof.}\ }}
\newtcolorbox{notebox}[1]{handoutbox, colback=hselight, colframe=hselight, before upper={\textcolor{hseblue}{\bfseries #1}\par\smallskip}}
\newtcolorbox{trapbox}{handoutbox, colback=hsewarm, colframe=hseaccent, leftrule=2.5pt, before upper={\textcolor{hseaccent}{\bfseries Trap.}\ }}
\newtcolorbox{selfcheckbox}{handoutbox, colback=white, colframe=hseaccent, leftrule=2.5pt, top=3pt, bottom=3pt, before upper={\textcolor{hseaccent}{\bfseries Check yourself.}\ }}
\newcommand{\handoutdisplay}[1]{\sbox0{$\displaystyle #1$}\ifdim\wd0>\linewidth
  \par{\centering\binoppenalty=10000\relpenalty=0\relax$\displaystyle #1$\par}\else\[#1\]\fi}
\newcommand{\selfcheckanswer}{\par\smallskip\textcolor{hsegray}{\bfseries Answer.}\ }

\setlist[itemize]{itemsep=3pt,topsep=4pt}
\setlist[enumerate]{itemsep=3pt,topsep=4pt}
\setlength{\parindent}{0pt}
\setlength{\parskip}{5pt plus 1pt}
\setlength{\emergencystretch}{2em}
\setcounter{tocdepth}{2}

\pagestyle{fancy}
\fancyhf{}
\fancyhead[L]{\small\color{hsegray}Optimization Methods}
\fancyhead[R]{\small\color{hsegray}Seminar 06}
\fancyfoot[L]{\small\color{hsegray}Evgeny Baulin}
\fancyfoot[C]{\small\thepage}
\fancyfoot[R]{\small\color{hsegray}MIT License}
\renewcommand{\headrulewidth}{0pt}
\renewcommand{\footrulewidth}{0pt}
\fancypagestyle{plain}{\pagestyle{fancy}}

\begin{document}

\begin{center}
  {\large\color{hsegray}Optimization Methods\par}
  \vspace{4pt}
  {\large\bfseries\color{hseaccent}Seminar 06\par}
  \vspace{8pt}
  {\LARGE\bfseries\color{hseblue}Second-Order Methods, Part Two: Gauss–Newton, Levenberg–Marquardt, Trust Region, L-BFGS\par}
  \vspace{8pt}
  {\Large Theory handout\par}
  \vspace{12pt}
  HSE University, Faculty of Computer Science\par
  Academic year 2026-2027\par
  \vspace{8pt}
  Course lecturer: Andrey Ignatov\par
  Evgeny Baulin\par
  Telegram: \href{https://t.me/tarakan_tuc}{@tarakan\_tuc}\par
\end{center}

\noindent{\color{hseblue}\rule{\linewidth}{0.6pt}}

\tableofcontents

\clearpage


Second-order methods, continued. Seminar 5 took the Newton step and its safeguards. This handout follows the lecture and the computer seminar of Andrey Ignatov on Gauss–Newton and Levenberg–Marquardt for least squares, the trust region as a different safeguard, and L-BFGS as the limited-memory quasi-Newton method.

\begin{conventionbox}
$r:\mathbb R^n\to\mathbb R^m$ is the residual with Jacobian $J$. $B$ is the current Hessian estimate, $G$ its inverse; $s_k=x_{k+1}-x_k$, $y_k=g_{k+1}-g_k$. The trust radius is $\Delta$; the quadratic model is $m(p)=f+g^\top p+\tfrac{1}{2}p^\top Bp$ and $\rho=(\text{actual})/(\text{predicted})$.
\end{conventionbox}

% =====================================================================
\section{Gauss–Newton}\label{sec:gn}

For $f(x)=\tfrac{1}{2}\|r(x)\|^2$ with $r:\mathbb R^n\to\mathbb R^m$ and Jacobian $J$, $\nabla f=J^\top r$ and $\nabla^2 f=J^\top J+\sum_i r_i\nabla^2 r_i$. Gauss–Newton keeps only $J^\top J$ and solves $J^\top Jp=-J^\top r$.

\begin{whybox}
Gauss–Newton is the Newton step for the linearized residual $r+Jp$. The dropped term is the one that would need second derivatives of $r$; it is small when the residuals at the solution are small.
\end{whybox}

\begin{theorembox}{Theorem: Convergence of Gauss–Newton.}
If $J(x^*)$ has full column rank and $r(x^*)=0$, there is a neighbourhood of $x^*$ in which Gauss–Newton converges to $x^*$ quadratically. If instead $\|r(x^*)\|$ is bounded away from zero, the convergence is at best linear.
\end{theorembox}

\begin{trapbox}
$J^\top J$ is positive semidefinite, not definite: a rank-deficient $J$ makes the normal equations singular. Levenberg–Marquardt below handles this with a shift; a QR factorization of $J$ is a more accurate alternative that never forms $J^\top J$.
\end{trapbox}

\problemhead{A1}{Gauss–Newton on an exponential fit}{On the board}
\textbf{Statement.} Fit $f_\theta(t)=a\,e^{bt}$ to the four points $(t_i,y_i)=\begin{pmatrix} 0 & 2 \\ 1 & 3.3 \\ 2 & 5.4 \\ 3 & 9 \end{pmatrix}$ with $\theta=(a,b)$. From $\theta_0=(1,\ 1)$ take two Gauss–Newton steps. Report $\theta_1$, $\theta_2$ and $\tfrac{1}{2}\|r(\theta_2)\|^2$.

\textbf{Solution.}
\begingroup\sloppy
\begin{enumerate}[label=\textbf{\arabic*.}, leftmargin=*]
  \item At $\theta_0$ the fitted values are $a_0e^{b_0t_i}=(1,\ 2.7183,\ 7.3891,\ 20.0855)$, so $r_0=(-1,\ -0.5817,\ 1.9891,\ 11.0855)$. The Jacobian has rows $(e^{b_0t_i},a_0t_ie^{b_0t_i})$: $J_0=\begin{pmatrix} 1 & 0 \\ 2.7183 & 2.7183 \\ 7.3891 & 14.7781 \\ 20.0855 & 60.2566 \end{pmatrix}$.
  \item The normal equations $J_0^\top J_0\,p=-J_0^\top r_0$ have $J_0^\top J_0=\begin{pmatrix} 466.416 & 1326.8717 \\ 1326.8717 & 3856.6408 \end{pmatrix}$ and $J_0^\top r_0=(234.7749,\ 695.7901)$. Solving by Cramer's rule gives $p_0=(0.4654,\ -0.3405)$, so $\theta_1=\theta_0+p_0=(1.4654,\ 0.6595)$.
  \item Repeat at $\theta_1$: predictions $(1.4654,\ 2.8337,\ 5.4798,\ 10.5967)$, residuals $(-0.5346,\ -0.4663,\ 0.0798,\ 1.5967)$, $J_1=\begin{pmatrix} 1 & 0 \\ 1.9338 & 2.8337 \\ 3.7395 & 10.9595 \\ 7.2314 & 31.79 \end{pmatrix}$. The normal equations give $p_1=(0.4557,\ -0.1548)$ and $\theta_2=(1.9211,\ 0.5047)$. The sum of squares falls from $64.0919$ to $1.5295$ to $0.0542$: by the second step we are already near the minimum.
  \item The exact second term of the Hessian is $\sum_i r_i\nabla^2 r_i$: for small residuals it is small, and Gauss–Newton behaves like Newton. The points sit on an exponential curve, so the residuals shrink fast and the convergence is almost quadratic.
\end{enumerate}
\endgroup
\begin{answerbox}
$\theta_1=(1.4654,\ 0.6595)$, $\theta_2=(1.9211,\ 0.5047)$, $\tfrac{1}{2}\|r(\theta_2)\|^2=0.0542$.
\end{answerbox}

% =====================================================================
\section{Levenberg–Marquardt}\label{sec:lm}

\begin{definitionbox}{Definition.}
\textbf{Levenberg–Marquardt} solves $(J^\top J+\lambda I)p=-J^\top r$ with $\lambda\ge0$ adapted by the ratio $\rho=\frac{f(x)-f(x+p)}{m(0)-m(p)}$: shrink $\lambda$ when $\rho$ is good, grow it when $\rho$ is bad.
\end{definitionbox}

Small $\lambda$ is Gauss–Newton; large $\lambda$ is a short gradient step $-J^\top r/\lambda$. The method interpolates between them without the user having to pick a line search strategy.

\problemhead{A2}{Levenberg–Marquardt rescues a wild start}{On your own}
\textbf{Statement.} On the fit of A1, start at the far point $\theta_0=(10,\ 1.5)$ and run three Levenberg–Marquardt steps with $\lambda_0=1$. The rule: compute $\rho=(f(\theta)-f(\theta+p))/(m(0)-m(p))$; if $\rho\ge0.25$ accept the step and set $\lambda\leftarrow\lambda/10$; otherwise reject and set $\lambda\leftarrow\lambda\cdot10$. Report $\theta$ after the first accepted step, after the third accepted step, and the number of rejections along the way.

\textbf{Solution.}
\begingroup\sloppy
\begin{enumerate}[label=\textbf{\arabic*.}, leftmargin=*]
  \item From $\theta_0$ the full Gauss–Newton step overshoots; LM with $\lambda_0=1$ shrinks it. The first few trials have $\rho<0.25$ and are rejected; $\lambda$ grows by ten each time until the shift is large enough.
  \item After the first accepted step $\theta=(0.8765,\ 1.474)$. Three accepted steps bring $\theta$ to $(1.0785,\ 0.8463)$; the number of rejections along the way is $0$.
  \item Each acceptance divides $\lambda$ by ten, each rejection multiplies it by ten. The method starts gradient-like, builds up Gauss–Newton behaviour as the model becomes trusted, and ends very close to the Gauss–Newton limit.
\end{enumerate}
\endgroup
\begin{answerbox}
After one accepted step $a=0.8765$; after three, $b=0.8463$; $0$ rejections.
\end{answerbox}

% =====================================================================
\section{Trust region}\label{sec:tr}

\begin{definitionbox}{Definition.}
A \textbf{trust-region step} minimizes the quadratic model over $\|p\|\le\Delta$:
\handoutdisplay{
p=\arg\min_{\|p\|\le\Delta} f+g^\top p+\tfrac{1}{2}p^\top Bp.
}
 By KKT, the minimizer satisfies $(B+\lambda I)p=-g$ for some $\lambda\ge0$ with $\lambda(\|p\|-\Delta)=0$.
\end{definitionbox}

The \textbf{dogleg} approximates the solution with the broken line from $0$ to the Cauchy point $p_C=-\alpha g$ to the Newton point $p_N=-B^{-1}g$: take the farthest point on it still inside the ball. It is exact when $B\succ0$ and the dogleg path is the trust-region path.

\begin{itemize}
  \item If $\rho>3/4$ and $\|p\|=\Delta$, double $\Delta$.
  \item If $\rho>1/4$, keep $\Delta$.
  \item If $\rho<1/4$, divide $\Delta$ by four.
  \item Reject the step whenever $\rho$ is too small; otherwise accept.
\end{itemize}

\problemhead{B1}{Dogleg on a stretched quadratic}{On the board}
\textbf{Statement.} On $f(x,y)=\tfrac{1}{2}(x^2+10\,y^2)$ at $x_0=(2,\ 1)$ with $B_0=\nabla^2 f$ and $\Delta=2$, compute the Cauchy point $p_C$ and the Newton point $p_N$. The Newton point sits outside the ball, so the dogleg follows the broken line from $0$ through $p_C$ to $p_N$ until it crosses the sphere. Give the second entry of $p_C$, the first entry of the dogleg step, and $f$ after the step.

\textbf{Solution.}
\begingroup\sloppy
\begin{enumerate}[label=\textbf{\arabic*.}, leftmargin=*]
  \item $g=Bx_0=(2,\ 10)$, $B=\operatorname{diag}(1,\ 10)$. The Newton step is $p_N=-B^{-1}g=(-2,\ -1)$ with norm $\|p_N\|=2.2361>\Delta$. The unconstrained minimizer is already outside the ball, so the dogleg will stop at the boundary.
  \item The quadratic $m(-\alpha g)=f-\alpha g^\top g+\tfrac{1}{2}\alpha^2 g^\top Bg$ is minimized at $\alpha^*=g^\top g/(g^\top Bg)=0.1036$; the resulting step has length $\alpha^*\|g\|=1.0564<\Delta$, so $p_C=-\alpha^*g=(-0.2072,\ -1.0359)$. Both coordinates were cut: the steep one more than the flat one.
  \item On the segment from $p_C$ to $p_N$, $p(\tau)=p_C+\tau(p_N-p_C)$, solve $\|p(\tau)\|^2=\Delta^2$ for the positive root: $\tau=0.8488$. The dogleg step is $p=(-1.7289,\ -1.0054)$ with $\|p\|=2=\Delta$, on the boundary.
  \item $f(x_0)=7$, $f(x_0+p)=0.0369$, model drop $m(0)-m(p)=6.9631$, so $\rho=1>3/4$. The step hit the boundary, so $\Delta$ grows; the next trial region has radius $4$.
\end{enumerate}
\endgroup
\begin{answerbox}
$p_C=(-0.2072,\ -1.0359)$, dogleg step $p=(-1.7289,\ -1.0054)$, $f=0.0369$.
\end{answerbox}

\problemhead{B2}{Trust region shrinks when the model lies}{On your own}
\textbf{Statement.} On $f(x)=-x+\tfrac{1}{4}x^4$ at $x_0=0.5$ with $B_0=f''(x_0)$, start with $\Delta_0=2$ and run the trust-region method until a step is accepted. The rule: reject when $\rho<0.25$ and set $\Delta\leftarrow\Delta/2$; accept otherwise. Give the accepted $\Delta$, the next iterate, and the number of rejections.

\textbf{Solution.}
\begingroup\sloppy
\begin{enumerate}[label=\textbf{\arabic*.}, leftmargin=*]
  \item At $x_0=0.5$: $g=-0.875$, $B=0.75$. The unconstrained minimizer of the model is $-g/B=1.1667$: it lies well outside the first ball of radius $2$.
  \item The first trial step has length $2$ and goes into the direction $-\operatorname{sign}(g)$. The quartic rises so fast that $f(x_0+p)>f(x_0)$ or $\rho<0.25$; $\Delta$ is divided by $2$ and the step tries again.
  \item The accepted $\Delta=0.5$ comes after $2$ rejections. The step is $0.5$ and $x_1=1$; $f$ falls from $-0.4844$ to $-0.75$.
\end{enumerate}
\endgroup
\begin{answerbox}
$\Delta=0.5$, $x_1=1$; $2$ rejections.
\end{answerbox}

% =====================================================================
\section{Limited memory BFGS}\label{sec:lbfgs}

Dense BFGS stores an $n\times n$ matrix. \textbf{L-BFGS} stores the last $m$ pairs $(s_i,y_i)$ and applies $G_k$ to $g_k$ without forming it: with scaling $\gamma_k=s_{k-1}^\top y_{k-1}/(y_{k-1}^\top y_{k-1})$ and the two-loop recursion, one direction costs about $4mn$ multiply-adds.

\begin{itemize}
  \item Set $q=g_k$. For the pairs from the newest to the oldest: $a_i=s_i^\top q/(y_i^\top s_i)$, $q\leftarrow q-a_iy_i$.
  \item Scale: $r=\gamma_kq$.
  \item For the pairs from the oldest to the newest: $b_i=y_i^\top r/(y_i^\top s_i)$, $r\leftarrow r+s_i(a_i-b_i)$.
  \item Return $p_k=-r$: about $4mn$ operations and $2mn$ stored numbers.
\end{itemize}

\begin{theorembox}{Proposition.}
The result equals $-Gg$ for the matrix $G$ obtained from $\gamma_kI$ by $m$ BFGS updates with the stored pairs. With $m\ge n$, L-BFGS and dense BFGS coincide to a floating-point tolerance.
\end{theorembox}

\problemhead{C1}{BFGS by hand}{On the board}
\textbf{Statement.} Let $f(x)=\tfrac{1}{2}x^\top Qx$ with $Q=\operatorname{diag}\begin{pmatrix} 1 & 0 \\ 0 & 4 \end{pmatrix}$, $x_0=(2,\ 2)$, $G_0=I$, and use exact line searches. Take two steps: form $s$, $y$ and the inverse BFGS update $G_1$ after the first step, check the secant condition, then take the second step. Report $G_1$ along the diagonal, $\alpha_1$ and $x_2$.

\textbf{Solution.}
\begingroup\sloppy
\begin{enumerate}[label=\textbf{\arabic*.}, leftmargin=*]
  \item $g_0=Qx_0=(2,\ 8)$, $p_0=-g_0$, $\alpha_0=g_0^\top g_0/(g_0^\top Qg_0)=0.2615$. So $s=\alpha_0p_0=(-0.5231,\ -2.0923)$, $x_1=(1.4769,\ -0.0923)$, $g_1=Qx_1=(1.4769,\ -0.3692)$ and $y=g_1-g_0=(-0.5231,\ -8.3692)$. On a quadratic $y=Qs$ exactly.
  \item $y^\top s=17.7846>0$, $\rho=1/(y^\top s)=0.0562$. With $V=I-\rho sy^\top$,
\handoutdisplay{
G_1=VG_0V^\top+\rho ss^\top=\begin{pmatrix} 1.0454 & -0.0028 \\ -0.0028 & 0.2502 \end{pmatrix}.
}

  \item Checks: $G_1y=(-0.5231,\ -2.0923)=s$, $G_1^\top=G_1$, eigenvalues $(0.2502,\ 1.0455)$, both positive. The inverse $B_1=\begin{pmatrix} 0.9566 & 0.0109 \\ 0.0109 & 3.9973 \end{pmatrix}$ differs from $Q$ in the direction orthogonal to $s$: one secant pair cannot pin down the whole matrix.
  \item $p_1=-G_1g_1=(-1.5451,\ 0.0966)$, $g_1^\top p_1=-2.3176<0$, exact step $\alpha_1=0.9559$: $x_2=(0,\ 0)=x^*=0$. With exact line searches BFGS minimizes an $n$-dimensional positive definite quadratic in at most $n$ steps.
\end{enumerate}
\endgroup
\begin{answerbox}
$G_1=\begin{pmatrix} 1.0454 & -0.0028 \\ -0.0028 & 0.2502 \end{pmatrix}$, $\alpha_1=0.9559$, $x_2=(0,\ 0)$.
\end{answerbox}

\problemhead{C2}{L-BFGS: two-loop recursion by hand}{On your own}
\textbf{Statement.} On the quadratic of C1 keep the last $m=2$ pairs. Take two steps of L-BFGS from $x_0=(2,\ 2)$ with exact line searches, apply the two-loop recursion with scaling $\gamma_k=s_{k-1}^\top y_{k-1}/y_{k-1}^\top y_{k-1}$. Report $\gamma_1$, the first entry of the direction $p_1=-H_1g_1$ computed by the recursion, and $x_2$.

\textbf{Solution.}
\begingroup\sloppy
\begin{enumerate}[label=\textbf{\arabic*.}, leftmargin=*]
  \item Step 0 is $-g_0$ with the exact step: it coincides with C1. After it, $\gamma_1=s_0^\top y_0/(y_0^\top y_0)=0.2529$.
  \item The two-loop pass with pairs $[(s_0,y_0)]$ and kernel $\gamma_1 I$ gives $p_1=(-0.3908,\ 0.0244)$. The exact step is $\alpha_1=3.7794$, so $x_2=(0,\ 0)$. On a 2-D quadratic L-BFGS with $m=n$ matches dense BFGS to the digit.
  \item The two-loop never forms $G$: it computes $Gg$ in $O(mn)$ operations. For $m\ll n$ it stores a Taylor-truncated version of the inverse Hessian along the last few directions.
\end{enumerate}
\endgroup
\begin{answerbox}
$\gamma_1=0.2529$, first entry of $p_1=-0.3908$, first entry of $x_2=0$.
\end{answerbox}

\begin{trapbox}
Pick $m$ by the problem: too small restarts the method every $m$ steps and loses iterations; too large wastes work without gain. The common range is between five and twenty, and the two-loop recursion rewards keeping the pair with the newest $s^\top y$ above zero.
\end{trapbox}

\end{document}
